% !TeX program = lualatex
% Compilar desde esta carpeta: lualatex -output-directory=build -interaction=nonstopmode -halt-on-error Propuesta_revistas.tex
\documentclass[11pt,a4paper,oneside]{report}
\usepackage{fontspec}
\setmainfont{Arial}
\usepackage[spanish,es-nodecimaldot,es-noshorthands]{babel}
\usepackage{geometry}
\geometry{a4paper,left=3cm,top=3cm,right=2.5cm,bottom=2.5cm,headheight=14pt,headsep=0.6cm}
\usepackage{setspace}
\onehalfspacing
\usepackage{graphicx}
\usepackage{booktabs}
\usepackage{longtable}
\usepackage{tabularx}
\usepackage{array}
\usepackage{pdflscape}
\usepackage{ragged2e}
\usepackage{xcolor}
\usepackage{colortbl}
\usepackage{fancyhdr}
\usepackage{titlesec}
\usepackage[hidelinks]{hyperref}
\usepackage{bookmark}
\usepackage{enumitem}
\usepackage{microtype}
\usepackage{tikz}
\usetikzlibrary{arrows.meta,positioning}
\usepackage{pifont}
\definecolor{unired}{HTML}{8B1A1A}
\definecolor{darkblue}{HTML}{17365D}
\definecolor{lightblue}{HTML}{DCE6F1}
\definecolor{baseline}{HTML}{5A7D71}
\definecolor{active}{HTML}{E9B44C}
\definecolor{planned}{HTML}{DCE6F1}
\setlength{\parindent}{0pt}
\setlength{\parskip}{9pt}
\setlength{\emergencystretch}{3em}
\setlist{leftmargin=1.1cm,itemsep=0.35\baselineskip,topsep=0.35\baselineskip}
\titleformat{\chapter}[hang]{\normalfont\bfseries\fontsize{11}{13.2}\selectfont}{\MakeUppercase{\chaptername}\ \thechapter.}{0.75em}{\MakeUppercase}
\titleformat{name=\chapter,numberless}[block]{\normalfont\bfseries\fontsize{11}{13.2}\selectfont}{}{0pt}{\MakeUppercase}
\titleformat{\section}{\normalfont\bfseries\fontsize{11}{13.2}\selectfont}{\thesection}{0.75em}{\MakeUppercase}
\titleformat{\subsection}{\normalfont\bfseries\fontsize{11}{13.2}\selectfont}{\thesubsection}{0.75em}{\MakeUppercase}
\titlespacing*{\chapter}{0pt}{0pt}{0.75\baselineskip}
\fancypagestyle{plain}{\fancyhf{}\fancyhead[R]{\thepage}\renewcommand{\headrulewidth}{0pt}}
\pagestyle{plain}
\newcommand{\autores}{LUIS ENRIQUE KOC G\'ONGORA\\HERBERT ANTONIO MEL\'ENDEZ GARC\'IA}
\newcommand{\asesor}{Samuel Oporto}
\newcommand{\caratula}[1]{%
  \hypersetup{pageanchor=false}
  \begin{titlepage}
  \thispagestyle{empty}
  \begin{center}
  {\bfseries\fontsize{19}{23}\selectfont UNIVERSIDAD NACIONAL DE INGENIER\'IA\par}
  \vspace{0.2cm}
  {\bfseries\fontsize{13}{16}\selectfont UNIDAD DE POSGRADO DE LA FACULTAD DE INGENIER\'IA INDUSTRIAL Y DE SISTEMAS\par}
  {\bfseries\fontsize{13}{16}\selectfont ESCUELA DE POSGRADO\par}
  \vspace{0.25cm}
  \includegraphics[width=4.6cm]{../Plan/imagenes/logo_UNI.png}\par
  \vspace{0.2cm}
  {\bfseries\fontsize{14}{17}\selectfont #1\par}
  \vspace{0.35cm}
  {\bfseries\fontsize{12}{15}\selectfont DISE\~NO Y EVALUACI\'ON DE UNA RED NEURONAL INFORMADA POR F\'ISICA PARA ESTIMAR LA TEMPERATURA Y LA AMPACIDAD DE CABLES EL\'ECTRICOS ENTERRADOS EN ENTORNOS HETEROG\'ENEOS\par}
  \vfill
  {\fontsize{11}{14}\selectfont PARA OBTENER EL GRADO ACAD\'EMICO DE\\MAESTRO EN CIENCIAS CON MENCI\'ON EN INTELIGENCIA ARTIFICIAL\par}
  \vspace{0.35cm}
  {\bfseries\fontsize{11}{14}\selectfont PRESENTADO POR:\par}
  {\bfseries\fontsize{11}{14}\selectfont \autores\par}
  \vspace{0.3cm}
  {\fontsize{11}{14}\selectfont ASESOR: \asesor\par}
  \vfill
  {\fontsize{11}{14}\selectfont Lima, Per\'u\\2026\par}
  \end{center}
  \end{titlepage}
  \hypersetup{pageanchor=true}
  \pagenumbering{arabic}
  \tableofcontents
  \clearpage
}
\begin{document}
\caratula{PROPUESTA DE REVISTAS PARA EL ART\'ICULO DERIVADO}

\chapter*{Criterio de selecci\'on}
\addcontentsline{toc}{chapter}{Criterio de selecci\'on}
La preselecci\'on prioriza revistas que ya han publicado estudios de cables subterr\'aneos, ampacidad o temperatura con suelo/backfill heterog\'eneo. La contribuci\'on prevista combina un artefacto PINN, verificaci\'on num\'erica/anal\'itica y aplicaciones t\'ermicas; la elecci\'on definitiva depender\'a de que el manuscrito sustente una validaci\'on rigurosa y diferencie con claridad los casos de literatura, benchmarks y activos peruanos.

\small
\setlength{\tabcolsep}{3pt}
\begin{longtable}{p{2.7cm}p{4.8cm}p{6.5cm}}
\toprule
\textbf{Revista} & \textbf{Art\'iculo comparable documentado} & \textbf{Encaje y recomendaci\'on} \\
\midrule
\endfirsthead
\toprule
\textbf{Revista} & \textbf{Art\'iculo comparable documentado} & \textbf{Encaje y recomendaci\'on} \\
\midrule
\endhead
\textit{Engineering Science and Technology, an International Journal} & Al-Dulaimi et al. (2024), ``Adaptive FEM-BPNN model for predicting underground cable temperature considering varied soil composition'', 51, 101658. DOI: \url{https://doi.org/10.1016/j.jestch.2024.101658}. & Encaje tem\'atico/m\'etodo muy alto: el art\'iculo ya cruza cable subterr\'aneo, composici\'on variable del suelo, predicci\'on de temperatura y FEM con red neuronal. Recomendada como primera opci\'on si el manuscrito demuestra ventaja metodol\'ogica frente a regresi\'on/FEM y presenta validaci\'on reproducible. \\
\textit{Results in Engineering} & Arias Velasquez et al. (2023), ``Enhancing medium voltage underground circuit design: Assessing limitations, thermal influence, and accurate modelling'', 20, 101552. DOI: \url{https://doi.org/10.1016/j.rineng.2023.101552}. & Encaje aplicado muy alto: existe un antecedente directo de dise\~no t\'ermico y modelaci\'on de circuitos subterr\'aneos de media tensi\'on, adem\'as de una conexi\'on peruana de autores/tema. Recomendada como alternativa prioritaria si el art\'iculo enfatiza utilidad de ingenier\'ia y evidencia de instalaciones peruanas. \\
\textit{Energies} & Atoccsa et al. (2024), ``Optimization of ampacity in high-voltage underground cables with thermal backfill using dynamic PSO and adaptive strategies'', 17(5), 1023. DOI: \url{https://doi.org/10.3390/en17051023}. & Encaje alto en ampacidad, cables enterrados y backfill; ofrece antecedente de autores peruanos y un caso de 220 kV. Es una opci\'on pertinente si el art\'iculo aporta un conjunto de datos claro y comparaciones de rating adem\'as del modelo PINN. \\
\bottomrule
\end{longtable}

\chapter*{Orden sugerido y comprobaciones previas}
\addcontentsline{toc}{chapter}{Orden sugerido y comprobaciones previas}
\begin{enumerate}
  \item \textit{Engineering Science and Technology, an International Journal}: mejor coincidencia combinada entre cable subterr\'aneo, heterogeneidad del suelo y aprendizaje/modelaci\'on.
  \item \textit{Results in Engineering}: alternativa con correspondencia especialmente directa en dise\~no de circuitos enterrados y modelaci\'on t\'ermica; permite destacar la pertinencia regional sin presentar inventarios p\'ublicos como geometr\'ia completa.
  \item \textit{Energies}: alternativa tem\'atica adecuada si el manuscrito pone en primer plano ampacidad, backfill y aplicaci\'on energ\'etica.
\end{enumerate}

\chapter*{Referencias verificadas}
\addcontentsline{toc}{chapter}{Referencias verificadas}
Al-Dulaimi, A. A., Guneser, M. T., Hameed, A. A., Garc\'ia M\'arquez, F. P., \& Gouda, O. E. (2024). Adaptive FEM-BPNN model for predicting underground cable temperature considering varied soil composition. \textit{Engineering Science and Technology, an International Journal, 51}, 101658. https://doi.org/10.1016/j.jestch.2024.101658\\
Arias Velasquez, R. M., Gonzales-Condori, E. G., \& Velarde Loayza, P. E. (2023). Enhancing medium voltage underground circuit design: Assessing limitations, thermal influence, and accurate modelling. \textit{Results in Engineering, 20}, 101552. https://doi.org/10.1016/j.rineng.2023.101552\\
Atoccsa, B. A., Puma, D. W., Mendoza, D., Urday, E., Ronceros, C., \& Palma, M. T. (2024). Optimization of ampacity in high-voltage underground cables with thermal backfill using dynamic PSO and adaptive strategies. \textit{Energies, 17}(5), 1023. https://doi.org/10.3390/en17051023
\end{document}